\documentclass[11pt]{article}
\usepackage[spanish]{babel}
\usepackage[utf8]{inputenc}
\usepackage[T1]{fontenc}
\usepackage{geometry}
\usepackage{graphicx}
\usepackage{booktabs}
\usepackage{amsmath}
\usepackage{hyperref}
\geometry{margin=2.5cm}

\newcommand{\maybeincludegraphics}[2][]{%
  \IfFileExists{#2}{\includegraphics[#1]{#2}}{\fbox{\parbox{0.88\linewidth}{Figura no disponible: \texttt{#2}}}}%
}

\title{Topological Incompleteness Index}
\author{Proyecto Mapper/TDA de exoplanetas}
\date{\today}

\begin{document}
\maketitle

\section{Introduccion}
Este subproyecto formaliza un indice topologico de incompletitud observacional para priorizar regiones Mapper y planetas ancla bajo lenguaje prudente.

\section{Motivacion}
El indice conecta sesgo de descubrimiento, sombra observacional, deficit local y planetas ancla como una extension metodologica del pipeline existente.

\section{Definicion de R3}
\[
x_i = (\log_{10} M_p, \log_{10} P, \log_{10} a).
\]

\section{Indice regional TOI}
\[
\mathrm{TOI}(v)=\mathrm{Shadow}(v)\,(1-I_{R^3}(v))\,C_{\mathrm{phys}}(v)\,S_{\mathrm{net}}(v).
\]

\section{Indice de ancla ATI}
\[
\mathrm{ATI}(p^*)=\mathrm{TOI}(v)\,\Delta_{\mathrm{rel,best}}(p^*)\,(1-I_{R^3}(p^*))\,A(p^*).
\]

\section{Deficit relativo de vecinos}
Se usa
\[
N_{\mathrm{obs}}(p^*,r), \qquad N_{\mathrm{exp}}(p^*,r), \qquad
\Delta_{\mathrm{rel}}(p^*,r)=\frac{N_{\mathrm{exp}}-N_{\mathrm{obs}}}{N_{\mathrm{exp}}+\epsilon}.
\]
El valor $\Delta_{\mathrm{rel,best}}$ ayuda a priorizar, pero debe interpretarse con cautela junto con el detalle por radio.

\section{Resultados}
\begin{figure}[ht]
\centering
\maybeincludegraphics[width=0.82\linewidth]{../../outputs/topological_incompleteness_index/figures_pdf/top_regions_toi_score.pdf}
\caption{Top 20 regiones Mapper por TOI.}
\end{figure}

\begin{figure}[ht]
\centering
\maybeincludegraphics[width=0.82\linewidth]{../../outputs/topological_incompleteness_index/figures_pdf/top_anchor_ati_score.pdf}
\caption{Top 20 planetas ancla por ATI.}
\end{figure}

\begin{figure}[ht]
\centering
\maybeincludegraphics[width=0.72\linewidth]{../../outputs/topological_incompleteness_index/figures_pdf/toi_vs_shadow_score.pdf}
\caption{Relacion entre shadow score y TOI.}
\end{figure}

\section{Top regiones}
Las regiones con TOI alto representan candidatas a submuestreo topologico bajo continuidad fisica, baja imputacion en $R^3$ y soporte de red suficiente.

\section{Top planetas ancla}
Los anclajes con ATI alto sirven para priorizacion observacional local, no como evidencia cerrada de objetos ausentes.

\section{Discusion}
Si se puede afirmar que TOI y ATI ordenan regiones y anclas donde el catalogo parece submuestreado bajo una referencia topologica. No se puede afirmar una cantidad absoluta de objetos reales ausentes.

\section{Limitaciones}
No hay funcion de completitud instrumental, $N_{\mathrm{exp}}$ es una referencia local, $\Delta_{\mathrm{rel,best}}$ puede inflar la lectura, $R^3$ simplifica la fisica, el proxy RV no es amplitud RV real, la imputacion puede afectar masa y los nodos pequenos pueden inflar indices.

\section{Conclusion}
TOI y ATI no detectan planetas faltantes, pero priorizan regiones y planetas ancla donde el catalogo parece topologicamente submuestreado.

\section{Trabajo futuro}
Completitud instrumental, inyeccion-recuperacion, validacion con catalogos futuros, integracion de propiedades estelares y comparacion entre metodos de descubrimiento.

\end{document}
